\documentclass[10pt,a4paper]{amsart}
\usepackage[margin=19mm]{geometry}
\usepackage{amsmath,amssymb,mathtools}
\usepackage[scr=rsfs,cal=euler]{mathalfa}
\usepackage{xfrac}
\usepackage[unicode,hidelinks,bookmarks=false]{hyperref}
\newcommand{\ie}{i.e.\ }
\newcommand{\cf}{cf.\ }
\newcommand{\resp}{respectively\ }
\newcommand{\Subsection}{\subsection}
\providecommand{\nobreakdash}{\nobreak\hbox{-}}
\setlength{\emergencystretch}{2em}
\multlinegap0pt
\usepackage{mdframed}
\usepackage{mdframed}
\usepackage{mdframed}
\usepackage{mdframed}
\usepackage{amssymb}
\usepackage{mdframed}
\usepackage{tikz}
\usepackage{mathtools}
\newtheorem{lem}{Lemma}
\renewcommand{\emptyset}{\varnothing}
\title{Finding Short Paths On Simple Polytopes}
\author{Alexander E. Black, Raphael Steiner}
\date{Source: arXiv:2603.05482v2 (CC BY 4.0)}
\begin{document}
\maketitle

\noindent\emph{Source and licence.} Finding Short Paths On Simple Polytopes. Version arXiv:2603.05482v2 (CC BY 4.0), distributed by arXiv under the Creative Commons Attribution 4.0 International licence, http://creativecommons.org/licenses/by/4.0/, as recorded in the arXiv licence metadata for this exact version and shown on the abstract page https://arxiv.org/abs/2603.05482v2. Full licence notice: \texttt{licenses/arxiv-2603.05482-CC-BY-4.0.txt}.\par\smallskip

\noindent\textit{Editorial scope.} Counted unit: the article's lem lem:partitionreduction (source0.tex lines 407--415). The article's complete original proof of it is delivered with it (source0.tex lines 417--467, one \texttt{proof} environment of 51 lines). No mathematics is written, repaired or re-derived: the statement and the proof are the article's own lines, in the article's own notation, and no retained line is substituted or re-flowed.  Retained uncounted as the source-local context the proof needs: lines 325--332; lines 379--390.  Nothing was omitted from any retained span, so the package carries no omission record.  No retained line contains a citation, so the wrapper carries no bibliography and no bibliography entry.  No retained line contains a figure environment, an image-inclusion command or drawing code, so no image is delivered and the package declares no asset dependency.  Editorial formatting only, in addition: an AMS \texttt{amsart} preamble that restates the article's own declarations and macros used by the retained lines, namely the environments \texttt{lem} on separate counters, together with the standard packages those lines need.  The retained environments print in this document as Lemma 1 in order of appearance, whereas the article prints the same items with the numbering of its own full document.  The complete native source of the article as distributed in the arXiv e-print for this exact version is bundled unchanged as \texttt{sources/195807/source0.tex}.
\begin{lem}
\label{lem:knapsackvertices}
The graph of $P_{\mathbf{b}}$ has vertex set $V_{1} \cup V_{2}$, where 
\begin{align*}
    V_{1} &= \left\{\mathbf{e}_S\bigg\vert S\subseteq [d+2] \text{ s. t.} \sum_{i \in S} w_{i} \leq \beta\right\} \text{ and} \\
    V_{2} &= \left\{\mathbf{e}_{S} + \frac{\beta+1/4-\sum_{i\in S} w_{i}}{w_{k}} \mathbf{e}_{k}\bigg\vert \sum_{i \in S} w_{i} < \beta + 1/4< \sum_{j \in S \cup \{k\}} w_{j} \text{ or }\sum_{i \in S} w_{i} > \beta + 1/4 > \sum_{j \in S \cup \{k\}} w_{j}\right\}.
\end{align*}
\end{lem}

\begin{lem}
\label{lem:edgetypes}
Let $\mathbf{b}\in \mathbb{Z}_{>0}^d$ be such that $\sum_{i=1}^d b_i$ is even and let $\mathbf{u}$ and $\mathbf{v}$ be two vertices of $P_{\mathbf{b}}$. Then $\mathbf{u}$ and $\mathbf{v}$ are adjacent on $P_\mathbf{b}$ if and only if 
\begin{itemize}
    \item[(a)] $\mathbf{u} = S$ and $\mathbf{v} = T$ for some $S, T\subseteq [d+2]$ with $|S\Delta T|=1$, or
    \item[(b)]  $\mathbf{u} = S$ and $\mathbf{v} = (S,i)$ (for $i\notin S$) or $\mathbf{v} = (S\setminus \{j\}, j)$ (for $j\in S$) for some $S\subseteq [d+2]$, or
    \item[(c)] $\mathbf{u} = (S,i)$ and $\mathbf{v} = (S,j)$ for some $S\subseteq [d+2]$ and distinct $i,j\notin S$, or
    \item[(d)] $\mathbf{u} = (S,i)$ and $\mathbf{v} = (S \cup \{i\}, j)$ for some $S\subseteq [d+2]$ and distinct $i,j\notin S$, or
    \item[(e)] $\mathbf{u} = (S,i)$ and $\mathbf{v} = (S \cup \{j\}, i)$ for some $S\subseteq [d+2]$ and distinct $i,j\notin S$, or
    \item[(f)] $\mathbf{u} = (S \cup \{i\},j)$ and $\mathbf{v} = (S \cup \{j\}, i)$ for some $S\subseteq [d+2]$ and distinct $i,j\notin S$.
\end{itemize}
\end{lem}

\begin{lem}
\label{lem:partitionreduction}
Let $\mathbf{b}=(b_1,\ldots,b_d)\in \mathbb{Z}_ {>0}^d$ such that $\sum_{i=1}^{d}{b_i}$ is even. Then
\begin{itemize}
    \item $(\emptyset,d+2)$ and $([d+1],d+2)$ are vertices of $P_{\mathbf{b}}$. 
    \item The shortest path between $(\emptyset,d+2)$ and $([d+1],d+2)$ is of length at most $d+1$ if and only if there exists a solution to \textsc{Partition  with even sum} with instance $\mathbf{b}$.
\end{itemize}

\end{lem}

\begin{proof}
Since 
\[\mathbf{w}^{\intercal} \mathbf{e}_{\emptyset} = 0 \leq \beta +1/4 < \beta +1/2 = \mathbf{w}^{\intercal} \mathbf{e}_{d+2},\] and 
\[\mathbf{w}^{\intercal} \mathbf{e}_{[d+1]} = \sum_{i=1}^{n} b_{i} -\beta = 2\beta - \beta = \beta < \beta + 1/4 < \beta + \beta +1/2 = \mathbf{w}^{\intercal} \mathbf{e}_{[d+2]},\]
$(\emptyset,d+2)$ and $([d+1],d+2)$ are vertices of $P_{\mathbf{b}}$ by the characterization of the vertices in Lemma \ref{lem:knapsackvertices}. This proves the first item of the lemma.

In the following, we will determine precisely the structure of the paths of length at most $d+1$ from $(\emptyset,d+2)$ to $([d+1],d+2)$ on $P_\mathbf{b}$, which will then yield the second item of the lemma.

From the characterization of edges in Lemma \ref{lem:edgetypes}, moving along any edge of $P_\mathbf{b}$ can only increase the size of the support of the current vertex by at most $1$. Thus, any path between $(\emptyset, d+2)$ and $([d+1], d+2)$ of length at most $d+1$ must in fact be of length \emph{exactly} $d+1$ and each step along the path must increase the size of the support by exactly $1$. The only edge types from Lemma \ref{lem:edgetypes} that increase the size of the support when we move along them starting from a vertex of the form $(S,i)$ are those of type (d) and (e). Since our path starts at $(\emptyset, d+2)$ and since moving along type (d) and (e) edges we stay within vertices of type $(S,i)$, it follows that any path of length $d+1$ from $(\emptyset,d+2)$ to $([d+1],d+2)$ on $P_\mathbf{b}$ must only use type (d) and (e) edges. We now claim that any such path in fact only uses type (e) edges. Indeed, towards a contradiction suppose it uses some type (d) edge and consider the earliest such edge along the path when starting from $(\emptyset, d+2)$. Since edges of type (e) always move from a vertex of the form $(S,i)$ to a vertex of the form $(S',i)$ and hence always preserve the ``second coordinate'', and since we start from the vertex $(\emptyset,d+2)$, the first edge of type (d) along the path must then start at a vertex of the form $(S,d+2)$ for some $S\subseteq [d+1]$ and go to $(S\cup \{d+2\},j)$ for some $j\notin S$ distinct from $d+2$. To have a total length of $d+1$, we would then need to reach $([d+1],d+2)$ from $(S\cup \{d+2\},j)$ using only type (d) and (e) edges which increase the support. However, this is impossible, since $S\cup \{d+2\}$ contains the element $d+2$ while $[d+1]$ does not, and since any type (d) and (e) edges used after will have to increase the support and hence preserve that $d+2$ is an element of the set in the tuple. Hence, we have reached the desired contradiction, and it follows that indeed all edges used along the path must be support-increasing edges of type (e).

Recalling the definition of type (e) edges, it now follows that every path of length at most $d+1$ from $(\emptyset, d+2)$ to $([d+1], d+2)$ on $P_\mathbf{b}$ must be of the form $(S_0,d+2),(S_1,d+2),\ldots,(S_{d+1},d+2)$ where \[\emptyset = S_{0} \subsetneq  S_{1} \subsetneq S_{2}, \dots \subsetneq S_{d+1} = [d+1]\] are such that $(S_{i},d+2)$ is a vertex of $P_\mathbf{b}$ and $S_{i} = S_{i-1} \cup \{k\}$ for some $k \in [d+1] \setminus S_{i-1}$ for all $1\le i \le d+1$.

We claim that such a sequence of sets exists (and hence the distance from $(\emptyset,d+2)$ to $([d+1],d+2)$ is at most $d+1$) if and only if there is a solution to \textsc{Partition with even sum}. Suppose first that such a sequence of sets exists. Let $i$ be minimal such that $d+1 \in S_{i}$. Then, since $(S_{i}, d+2)$ is a vertex of $P_{\mathbf{b}}$,

\begin{align*}
 -\beta +\sum_{j \in S_{i-1}} b_{j} &= \mathbf{w}^{\intercal}\mathbf{e}_{d+1} + \mathbf{w}^{\intercal} \mathbf{e}_{S_{i-1}}\\
 &= \mathbf{w}^{\intercal} \mathbf{e}_{S_{i}} \\
 &\leq \beta + 1/4 \\
 &\leq \mathbf{w}^{\intercal}\mathbf{e}_{S_{i} \cup \{d+2\}} \\
 &= \mathbf{w}^{\intercal} \mathbf{e}_{d+2} + \mathbf{w}^{\intercal} \mathbf{e}_{S_{i}} \\
 &=  \beta + 1/2 - \beta + \sum_{j \in S_{i-1}} b_{j} \\
 &= 1/2 + \sum_{j \in S_{i-1}} b_{j}.      
\end{align*}

In particular, $\beta + 1/4 \leq 1/2 + \sum_{j \in S_{i-1}} b_{j}$, so
\[\sum_{j \in S_{i-1}} b_{j} \geq \beta -1/4.\]

Similarly, since $(S_{i-1},d+2)$ is also a vertex of $P_{\mathbf{b}}$ and since $S_{i-1}$ does not contain $d+1$ by definition of $i$, we have:
\[\sum_{j \in S_{i-1}}b_{j} = \mathbf{w}^{\intercal}\mathbf{e}_{S_{i-1}} \leq \beta + 1/4.\]
It follows that 
\[\beta - 1/4 \leq \sum_{j \in S_{i-1}} b_{j} \leq \beta + 1/4\]
Since $b_{i} \in \mathbb{Z}$ for all $i \in [n]$ and $\beta\in \mathbb{Z}$, it follows that $\sum_{j \in S_{i-1}} b_{j} = \beta$. Therefore, in that case, \textsc{Partition with even sum} has a solution.

Suppose instead that \textsc{Partition with even sum} has a solution. Up to reordering we may without loss of generality assume then that
\[\sum_{i=1}^{k} b_{i} = \beta.\]
For each $j\in [d+1]$, define 
\[S_{j} =\begin{cases}
    \{1, \dots,j\} &\text{if } j \leq k \\ \\
    \{1,\dots, j-1\} \cup \{d+1\} &\text{if } j\ge k+1. 
\end{cases}\]
Then it suffices to show that $(S_{j}, d+2)$ is a vertex for each $j \in [d+1]$. If $j \leq k$, then 
\[\mathbf{w}^{\intercal} \mathbf{e}_{S_{j}} = \sum_{i \in S_{j}} b_{i} = \sum_{i=1}^{j} b_{i} \leq \sum_{i=1}^{k} b_{i} \leq \beta < \beta + 1/4 < \beta + 1/2 + \sum_{i\in S_{j}} b_{i} = \mathbf{w}^{\intercal} \mathbf{e}_{S_{j} \cup \{d+2\}}.\]
Hence, $(S_{j}, d+2)$ is a vertex in that case. 

If $j = k+1$, then
\[\mathbf{w}^{\intercal} \mathbf{e}_{S_{j}} = -\beta + \sum_{i=1}^{k} b_{i} = 0 < \beta + 1/4 < \beta + 1/2 + 0 = w_{d+2} + \sum_{i \in S_{j}} w_{i} = \mathbf{w}^{\intercal} \mathbf{e}_{S_{j} \cup \{d+2\}}.\]

Finally, suppose that $j \geq k+2$. Then $\sum_{i \in S_{j}} w_{i} \geq \sum_{i\in S_{k+1}}w_i=0$, and
\[\mathbf{w}^{\intercal}\mathbf{e}_{S_{j}} = \sum_{i \in S_{j}} w_{i} \leq \sum_{i \in [d+1]} w_{i} = \beta < \beta + 1/4 < \beta + 1/2 \leq \beta + 1/2 + \sum_{i \in S_{j}} w_{i} = \mathbf{w}^{\intercal} \mathbf{e}_{S_{j} \cup \{d+2\}}.\]
Hence, in all cases, $(S_{j}, d+2)$ is a vertex and so there is a path from $(\emptyset, d+2)$ to $([d+1], d+2)$ of length at most $d+1$ of the desired form. This concludes the proof of the equivalence claimed in the second item of the lemma. 
\end{proof}

\end{document}
